Ich erkläre, dass ich bei der Erstellung dieser Arbeit generative KI-Werkzeuge
verwendet habe. Sie kamen in drei voneinander unterscheidbaren Rollen zum
Einsatz; für jede ist im Folgenden angegeben, wo in der Arbeit sie zum Tragen
kommt.

\emph{Als Schreibhilfe.} Entwerfen, Umstrukturieren, Paraphrasieren sowie
Prüfung von Rechtschreibung, Grammatik und der Konsistenz von Terminologie und
Zahlenangaben über die Kapitel hinweg.

\emph{Als Entwicklungshilfe.} Codegenerierung, Refactoring, Erstellung von Tests
und Fehlersuche für das in Kapitel~\ref{chap:implementation} beschriebene
Software-Artefakt sowie für die \LaTeX{}-Quellen dieses Dokuments.

\emph{Als Forschungsinstrument.} Ein Vision-Language-Modell wurde eingesetzt, um
das bildweise Kodierschema auf die Bildschirmaufzeichnungen der Evaluierungsstudie
anzuwenden. Dieser Einsatz ist keine Schreibhilfe, sondern Teil der Analyse und
daher nicht nur hier, sondern in der Arbeit selbst dokumentiert: Das
Kodierprotokoll findet sich in Abschnitt~\ref{app:eval-stagelog}, die
Übereinstimmung zwischen der Kodierung des Modells und jener des Autors über die
von beiden abgedeckten Einzelbilder samt den beiden systematischen Quellen der
Abweichung in Abschnitt~\ref{app:eval-agreement}. Kein Ergebnis der Arbeit beruht allein auf der Kodierung des Modells.

Alle Ausgaben wurden kritisch geprüft und, wo angebracht, überarbeitet. Mein
gestalterischer Einfluss überwiegt in der vorliegenden Arbeit. Die
Entwurfsentscheidungen, die Interpretation der Evidenz und die gezogenen
Schlussfolgerungen sind meine eigenen. Verwendete KI-Werkzeuge:

\begin{itemize}
  \item Claude\footnote{\url{https://claude.ai}}, Anthropic. Modelle: Claude
        Sonnet~4.5 und Claude Opus~4.1--5, Versionsstand März~2026 --
        September~2026. Schreib- und Entwicklungshilfe.
  \item Gemini\footnote{\url{https://gemini.google.com}}, Google. Modell:
        Gemini~3.1~Pro, Versionsstand Oktober~2025 -- September~2026.
        Schreibhilfe.
  \item Otter.ai\footnote{\url{https://otter.ai}}. Automatische
        Spracherkennung, Versionsstand März~2026 -- Juli~2026. Transkription
        der Audioaufzeichnungen beider Studien
        (Kapitel~\ref{chap:evaluation}, Abschnitt~\ref{sec:eval-tworecords}).
  \item Qwen3-VL-30B-A3B-Instruct\footnote{\url{https://github.com/QwenLM/Qwen3-VL}},
        Alibaba Cloud, lokal ausgeführt. Versionsstand Juli~2026. Bildweise
        Kodierung der Bildschirmaufzeichnungen der Evaluierungsstudie
        (Abschnitt~\ref{app:eval-stagelog}).
\end{itemize}
